% ID: SEL-D0010-BC-E06
% ESTADO: validada. S011; tres pasadas del mismo autor, sin segundo revisor.
\textbf{a) Unión de los dos medios de transporte.}\par\nopagebreak[4]
Sean $V$ tener vehículo propio y $T$ usar transporte público. El dato de no usar ninguno es $P(V^c\cap T^c)=0.03$. Por De Morgan,
\[
\boxed{P(V\cup T)=1-0.03=0.97.}
\]
La unión es inclusiva: incluye a quienes cumplen ambas condiciones.

\textbf{b) Transporte público sin vehículo propio.}\par\nopagebreak[4]
Primero usamos inclusión-exclusión:
\[
P(V\cap T)=P(V)+P(T)-P(V\cup T)=0.90+0.40-0.97=0.33.
\]
Quitamos esa intersección del transporte público:
\[
\boxed{P(T\cap V^c)=0.40-0.33=0.07.}
\]
\textbf{c) Probabilidad condicionada.}\par\nopagebreak[4]
El espacio de referencia pasa a ser el de quienes no tienen vehículo, de probabilidad $0.10$:
\[
\boxed{P(T\mid V^c)=\frac{P(T\cap V^c)}{P(V^c)}=\frac{0.07}{0.10}=0.70.}
\]
Comprobación: las cuatro regiones disjuntas son $0.33$ (ambos), $0.57$ (sólo vehículo), $0.07$ (sólo transporte público) y $0.03$ (ninguno); suman $1$ y reproducen los dos marginales. No se ha supuesto independencia.
